\documentclass[10pt,a4paper]{amsart}
\usepackage[margin=19mm]{geometry}
\usepackage{amsmath,amssymb,mathtools}
\usepackage[scr=rsfs,cal=euler]{mathalfa}
\usepackage{xfrac}
\usepackage[unicode,hidelinks,bookmarks=false]{hyperref}
\newcommand{\ie}{i.e.\ }
\newcommand{\cf}{cf.\ }
\newcommand{\resp}{respectively\ }
\newcommand{\Subsection}{\subsection}
\providecommand{\nobreakdash}{\nobreak\hbox{-}}
\setlength{\emergencystretch}{2em}
\multlinegap0pt
\theoremstyle{plain}
\newtheorem{thm}{Theorem}
\newtheorem{assumption}[thm]{Assumption}
\newtheorem{lemma}[thm]{Lemma}
\newtheorem{prop}[thm]{Proposition}
\newtheorem{cor}[thm]{Corollary}
\theoremstyle{definition}
\newtheorem{definition}[thm]{Definition}
\theoremstyle{remark}
\newtheorem{rmk}[thm]{Remark}
\usepackage{braket}
\newcommand{\abs}[1]      {\left| #1 \right|}
\newcommand{\norm}[1]     {\left\| #1 \right\|}
\newcommand{\ot}{\otimes}
\newcommand{\lrb}[1]{\left ( #1 \right )}
\newcommand{\lrcb}[1]{\left \{ #1 \right \}}
\newcommand{\lrsb}[1]{\left [ #1 \right ]}
\newcommand{\diag}[1]{\operatorname{diag}\lrb{#1}}
\newcommand{\realpart}[1]{\operatorname{Re}\lrb{#1}}
\newcommand{\imaginarypart}[1]{\operatorname{Im}\lrb{#1}}
\newcommand{\mypoly}[1]{\operatorname{poly} \lrb{#1}}
\newcommand{\mylog}[1]{\operatorname{log} \lrb{#1}}
\newcommand{\myO}[1]{O\lrb{#1}}
\newcommand{\mytO}[1]{\tilde{O}\lrb{#1}}
\newcommand{\myOmega}[1]{\Omega\lrb{#1}}
\newcommand{\myTheta}[1]{\Theta\lrb{#1}}
\let\dotlessi\i
\renewcommand{\i}{\ensuremath{\mathbf{i}}}
\newcommand{\R}{\mathbb{R}}
\title[Quantum algorithms for general nonlinear dynamics based on the Carleman embedding]{Quantum algorithms for general nonlinear dynamics based on the Carleman embedding: BQP-completeness of Poly-CNOP (printed as Theorem 7.6)}
\author{David Jennings, Kamil Korzekwa, Matteo Lostaglio, Andrew T Sornborger, Yi\u{g}it Suba\c{s}\dotlessi, Guoming Wang}
\date{Source: arXiv:2509.07155v2 (CC BY 4.0)}
\begin{document}
\maketitle

\noindent\emph{Source and licence.} Quantum algorithms for general nonlinear dynamics based on the Carleman embedding. Version arXiv:2509.07155v2 (CC BY 4.0), distributed under the Creative Commons Attribution 4.0 International licence, http://creativecommons.org/licenses/by/4.0/, as recorded in the arXiv licence metadata for this exact version and shown on the saved abstract page https://arxiv.org/abs/2509.07155v2. Full licence notice: licenses/arxiv-2509.07155-CC-BY-4.0.txt.\par\smallskip

\noindent\textit{Editorial scope.} Counted unit: the BQP-completeness of the decision problem Poly-CNOP for coupled nonlinear oscillators, printed in the article as Theorem 7.6 (source0.tex lines 4113--4115, one \texttt{thm} environment), delivered together with the article's complete original proof of it (source0.tex lines 4117--4216, one \texttt{proof} environment of 100 lines immediately adjacent to the statement, with no gap and no deferral). No mathematics is written, repaired or re-derived: statement and proof are the article's own text line for line, in the article's own notation, and no line of either is omitted, substituted or re-flowed.

Required context retained uncounted: the article's definition of the Coupled Nonlinear Oscillator Problem (lines 4053--4077) and of Poly-CNOP (lines 4079--4110), which the statement names; the article's own definition of the graph Laplacian matrix with its labelled display (lines 3787--3791) and of the signed incidence matrix with its labelled display (lines 3874--3878), both of which the definition of Poly-CNOP refers to; the article's complexity theorem for time-dependent systems with a quadratic nonlinearity, whose algorithm the proof invokes five times (lines 3660--3689); and the article's own section and subsection headings that the retained cross-references point at (lines 3574, 3584--3585, 3763 and 3769--3770 and 3774--3775). Those pieces are exactly what the statement names and what the proof uses; the proof's remaining steps are its own.

References. Exactly three citation commands occur in the retained spans, all in the article's own sentences: the BQP-hardness reference, the reference for the incidence-matrix factorisation and the reference for the algorithm it follows. The package therefore carries exactly those three bibliography entries, transcribed verbatim from the article's own compiled bibliography and printed with the article's own labels 24, 44 and 45.

Editorial formatting only, in addition: an AMS \texttt{amsart} document that restates the macro definitions the retained lines use (the article's own \texttt{\textbackslash R}, \texttt{\textbackslash abs}, \texttt{\textbackslash norm}, \texttt{\textbackslash ot}, \texttt{\textbackslash lrb}, \texttt{\textbackslash lrcb}, \texttt{\textbackslash lrsb}, \texttt{\textbackslash diag}, \texttt{\textbackslash realpart}, \texttt{\textbackslash imaginarypart}, \texttt{\textbackslash mypoly}, \texttt{\textbackslash mylog}, \texttt{\textbackslash myO}, \texttt{\textbackslash mytO}, \texttt{\textbackslash myOmega}, \texttt{\textbackslash myTheta} and \texttt{\textbackslash i}), loads the standard braket package for the article's ket and bra notation, and restates the article's own theorem-environment declarations. Consequently the counted theorem prints here as Theorem 4, whereas the article prints it as Theorem 7.6, and the retained environments and equations are numbered anew in order of appearance, so the article's own numbers are not reproduced; every cross-reference inside the retained text resolves to this wrapper's numbering, and that is the only change to the numbering. No figure, no image-inclusion line and no drawing code occur in any retained line, so no artwork is delivered or needed, although the article contains figures elsewhere.

The article's publisher class file \texttt{quantumarticle.cls} and its bibliography style file and database are not shipped in this package and are not used to build it. Not retained: the article's other results, its appendices, its figures and its bibliography except the three cited entries; none of that material is used as a step of the delivered proof, which is complete as the author wrote it. The package ships the article's concatenated source verbatim as \texttt{sources/184703/source0.tex}, and every delivered excerpt is an exact line range of that file. No mathematical statement, hypothesis, estimate or conclusion has been rewritten, repaired or added, and printed slips, if any, are preserved literally.
\section{Quantum algorithms for nonlinear dynamical systems}

\subsection{Nonresonant and certain resonant nonlinear problems}
\label{subsec:quantum_algorithm_nonresonant_systems}

\begin{thm}
Consider the $N$-dimensional quadratic ODE system
\begin{align}
    \dot{x}(t) = F_1 x(t) + e^{\i \omega t}F_2 x(t)^{\ot 2},
\end{align}
where $\omega \in \R^+$ \footnote{A similar result holds for $\omega \in \R^-$ as well. This extension is straightforward and hence we leave the details to the reader.}, 
$F_1=Q\Lambda Q^{-1}$ with $\Lambda=\diag{\lambda_1, \lambda_2, \dots, \lambda_N}$, $\realpart{\lambda_i} \le 0$ and $\abs{\imaginarypart{\lambda_i}} \le \omega/4$, for all $i \in [N]$. Let $\tilde{F}_2=Q^{-1}F_2Q^{\otimes 2}$, and assume that
\begin{align}
32\rho s \norm{Q^{-1}} \norm{\tilde{F}_2} \le c\omega,    
\end{align}
where $\rho \ge \max_{t \in [0, T]} \norm{x(t)}$, $s=\operatorname{cs}(\tilde{F}_2)$, and 
$c \in (0,1/2)$ is a constant.

Let $U_i$ be a $(\alpha_i, a_i, 0)$-block-encoding of $F_i$, for $i=1,2$, and let $V_0$ be a unitary that prepares the normalized initial state $\ket{x(0)}$. Then for any $\epsilon \in (0, 1)$, we can prepare a quantum state that is $\epsilon$-close (in Euclidean norm) to $\ket{x(T)}$ using 
$$\mytO{\alpha g T  2^k \kappa_Q^k\mylog{1/\epsilon}}$$
queries to (controlled-) $U_1$, $U_2$, $V_0$ and their inverses, where 
\begin{align}
g=\frac{\max_{t \in [0, T]} {\norm{x(t)}}}{\min_{t \in [0, T]} {\norm{x(t)}}},~
\kappa_Q=\norm{Q}\cdot \norm{Q^{-1}},~
\alpha=\alpha_1+{\rho\kappa_Q\alpha_2},~ 
\end{align} 
and
\begin{align}
k=\myO{\max(\mylog{ g  \kappa_Q T\omega/(s\epsilon)}, 1)}.    
\end{align}
In addition to these queries, the algorithm uses
$$\mytO{\alpha g T 2^k \kappa_Q^k \mylog{1/\epsilon}\cdot \mypoly{\mylog{N}}}$$
elementary quantum gates.
\label{thm:algorithm_complexity_time_dependent_f2_case}
\end{thm}

\section{Exponential quantum advantage for nonlinear problems}

\subsection{Coupled oscillator systems}
\label{subsec:coupled_oscillators}

\subsubsection{Linear oscillators}
\label{subsubsec:coupled_linear_oscillators}

The coupling structure of the system can be concisely described using graph theory. We represent the network of springs as a weighted graph $G$, where the vertex set is $[N]$ and the edge weights are given by $\lrcb{\kappa_{i,j}:~1\le i\le j\le N}$. Note that $G$ may contain self-loops, corresponding to springs that connect individual masses to fixed walls. The \emph{Laplacian} matrix of $G$ is defined as
\begin{align}
L_G = \sum_{i=1}^N \lrb{\sum_{j=1}^N \kappa_{i,j}} \ket{i}\bra{i} - \sum_{1\le i<j\le N} \kappa_{i,j} \lrb{\ket{i}\bra{j}+\ket{j}\bra{i}}.
\label{eq:laplacian_matrix_def}
\end{align}

The Laplacian matrix $L_G$ admits a convenient factorization into two rectangular matrices that are transposes of each other\footnote{This decomposition of $L_G$ has been leveraged to develop several algorithms, including the ones presented in Refs.~\cite{wang2017efficient} and \cite{chakraborty2019thepower}.}. Specifically, the \emph{signed incidence} matrix of $G$ is an $N \times \frac{N(N+1)}{2}$ matrix defined by
\begin{align}
B_G = \sum_{1\le i<j\le N} \sqrt{\kappa_{i,j}} \left( \ket{i} - \ket{j} \right) \bra{i,j} + \sum_{i=1}^N \sqrt{\kappa_{i,i}}\ket{i}\bra{i,i}.
\label{eq:incidence_matrix_def}
\end{align}

\begin{definition}[Coupled Nonlinear Oscillator Problem (CNOP)]
Let $M=\diag{m_1,m_2,\dots,m_N}$ be an $N \times N$ diagonal matrix of masses $m_i>0$, and let $K=(\kappa_{i,j})_{1\le i,j\le N}$ be the $N\times N$ matrix of spring constants $\kappa_{i,j} \ge 0$, assumed to be $d$-sparse. Consider a network of coupled nonlinear oscillators governed by the equation:
\begin{align}
M\ddot{x}(t) = -L_G x(t)+ e^{\i \omega t}\lrsb{Q_1 (\dot{x}(t) \otimes \dot{x}(t)) + Q_2 (\dot{x}(t) \otimes x(t)) + Q_3 (x(t)\otimes x(t))
},
\end{align}
where $\omega \in \R^+$ and $L_G$ is defined as in Eq.~\eqref{eq:laplacian_matrix_def}. Suppose
\begin{align}
Q_1 = \sqrt{M} G_1 \lrb{\sqrt{M}\otimes \sqrt{M}},
~Q_2 = -\sqrt{M} G_2 \lrb{\sqrt{M} \otimes B_G^\dagger}, 
~Q_3 = -\sqrt{M} G_3 \lrb{B_G^\dagger \otimes B_G^\dagger},   
\end{align}
where $B_G$ is defined as in Eq.~\eqref{eq:incidence_matrix_def}, and
$G_1 \in \R^{N \times N^2}$, $G_2 \in \R^{N \times \frac{N^2(N+1)}{2}}$,
$G_3 \in \R^{N \times \frac{N^2(N+1)^2}{4}}$ are arbitrary.
Define the normalized state 
\begin{align}
\ket{\psi(t)}=\frac{1}{\sqrt{2E(t)}}\begin{pmatrix}
    \sqrt{M}\dot{x}(t) \\
    \i B_G^\dagger x(t)
\end{pmatrix},
\end{align} 
where $E(t)=\frac{1}{2}\dot{x}(t)^\dagger M \dot{x}(t)+\frac{1}{2}x(t)^\dagger L_G x(t)$. Suppose we are given access to oracles that compute the entries of $K$ and $M$, block-encodings of $G_1$, $G_2$ and $G_3$, along with a unitary $\mathcal{W}$ that prepares $\ket{\psi(0)}$. 
Let $V \subseteq [N]$ be a subset for which there exists an oracle that decides its membership. Given a time $T>0$ and real numbers $0\le a < b \le 1$, the goal is to decide whether $K_V(T)/E(T) <a$ or $K_V(T)/E(T) > b$, promised that one of these holds.    
\end{definition}

\begin{definition}[Poly-CNOP]
The setup follows that of CNOP, with additional constraints on the input encodings and parameter settings. Specifically, we require that $N=2^n$,  
\begin{align}
d, T, \kappa_{\max}, m_{\max}, \frac{1}{m_{\min}}, E_{\rm max}, \frac{1}{E_{\rm min}} = \myO{\operatorname{poly}(n)}, 
\end{align}
where $m_{\max} = \max_{i \in [N]} m_i$, $m_{\min} = \min_{i \in [N]} m_i$, $\kappa_{\max} = \max_{i,j \in [N]} \kappa_{i,j}$, $E_{\rm max}=\max_{t \in [0, T]} E(t)$
and
$E_{\rm min}=\min_{t \in [0, T]} E(t)$, and
\begin{align}
b - a = \myOmega{\frac{1}{\operatorname{poly}(n)}}.    
\end{align}
 
Moreover, let $\rho=\myO{\mypoly{n}}$ be a known upper bound on $E_{\rm max}$. Suppose $\bar{B}_G$ has the singular value decomposition $\bar{B}_G=U \Gamma W^\dagger$, where $\Gamma=\sum_{j=1}^N \gamma_j \ket{j}\bra{j}$ is diagonal, and $U=\sum_{j=1}^N \ket{u_j}\bra{j}$ and $W=\sum_{j=1}^{N} \ket{v_j}\bra{j}$ are unitary \footnote{Technically, $\bar{B}_G$ is a rectangular matrix and $W$ is an isometry, as each $\ket{v_j}$ resides in a $\frac{N(N+1)}{2}$-dimensional space. However, the system's state always lies within the $N$-dimensional subspace spanned by $\ket{v_1}, \ket{v_2}, \dots, \ket{v_N}$. Therefore, we restrict our attention to this subspace, effectively treating $\bar{B}_G$ as a square matrix and $W$ as a unitary.}. Define
\begin{align}
\tilde{G}_1=U^\dagger G_1 (U \otimes U),~
\tilde{G}_2=U^\dagger G_2 (U \otimes W),~
\tilde{G}_3=U^\dagger G_3 (W \otimes W).
\end{align}
Let $s_i=\operatorname{cs}(\tilde{G}_i)$, for $i=1,2,3$ \footnote{As in Section~\ref{subsec:quantum_algorithm_nonresonant_systems}, these column-sparsity parameters enter through the convergence condition for the Carleman truncation. They should not be interpreted as additional cost parameters in the assumed block-encodings of $G_1$, $G_2$ and $G_3$.}. We assume that
\begin{align}
    64 \rho (s_1+s_2+s_3)\sqrt{\norm{\tilde{G}_1}^2+\norm{\tilde{G}_2}^2+\norm{\tilde{G}_3}^2}\le c \omega
\end{align}
where $c \in (0,1/2)$ is a constant. In addition, we assume that 
\begin{align}
\omega \ge 4\sqrt{2d\kappa_{\rm max}/m_{\rm min}}, \quad \omega = \myO{(s_1+s_2+s_3) \cdot\mypoly{n}}.
\end{align}


Furthermore, we assume that the oracles for $K$, $M$, $V$, and the unitary $\mathcal{W}$ that prepares $\ket{\psi(0)}$ can be implemented by quantum circuits of $\myO{\operatorname{poly}(n)}$ sizes. In addition, for each $i=1,2,3$, 
a block-encoding of $G_i$ with a rescaling factor $\myO{\mypoly{n}}$ can be implemented using $\myO{\mypoly{n}}$ elementary gates. 

\end{definition}

\begin{thm}
Poly-CNOP is BQP-complete.
\end{thm}

\begin{proof}
Since Poly-CNOP generalizes Poly-CLOP -- corresponding to the special case where $G_1$, $G_2$ and $G_3$ are all zero -- and the latter has been shown to be BQP-hard in Ref.~\cite{babbush2023exponential}, it follows that Poly-CNOP is also BQP-hard.

It remains to prove that Poly-CNOP can be solved in $\myO{\mypoly{n}}$ quantum time. This is accomplished by embedding the dynamics of the coupled nonlinear oscillators into a nonlinear Schr\"{o}dinger equation and applying the quantum algorithm from Theorem~\ref{thm:algorithm_complexity_time_dependent_f2_case} to solve it.

Specifically, as mentioned earlier, the vector
\begin{align}
z(t)=\begin{pmatrix}
    \sqrt{M}\dot{x}(t)\\
    \i B_G^\dagger x(t)
\end{pmatrix} 
\end{align}
evolves according to the differential equation
\begin{align}
\dot{z}(t) = -\i H z(t) + e^{\i \omega t} F_2 z(t)^{\otimes 2},    
\label{eq:nonlinear_schrodinger_equation_for_nonlinear_oscillators}
\end{align}
where
\begin{align}
H = -\begin{pmatrix}
0 & \bar{B}_G \\
\bar{B}_G^{\dagger} & 0
\end{pmatrix},~&~
F_2=\begin{pmatrix}
    G_1 & \i G_2 & 0 & G_3 \\
    0 & 0 & 0 & 0
\end{pmatrix}. 
\end{align}

Next, we show that this nonlinear Schr\"{o}dinger equation satisfies all the conditions required by Theorem~\ref{thm:algorithm_complexity_time_dependent_f2_case}:
\begin{itemize}
\item First, in this problem, we have $F_1=-\i H$, where $H$ is Hermitian and has
eigenvalues $\pm \gamma_1, \pm \gamma_2$, $\dots$, $\pm \gamma_N$. As shown in Section~\ref{subsubsec:coupled_linear_oscillators}, each $\gamma_j$ satisfies $\gamma_j\le \sqrt{2 d \kappa_{\rm max} / m_{\rm min}}$. Thus, $F_1$ is anti-Hermitian, and all eigenvalues of $F_1$ have zero real parts, and imaginary parts lying in the interval $[-\omega/4, \omega/4]$.
\item Second, $F_1=-\i H$ admits the spectral decomposition
\begin{align}
    F_1 = \i\sum_{j=1}^N \gamma_j \lrb{\ket{+_j}\bra{+_j} - \ket{-_j}\bra{-_j}},    
\end{align}
where the eigenvectors are given by
\begin{align}
    \ket{\pm_j}=\frac{1}{\sqrt{2}}\lrb{\ket{0}\ket{u_j}\pm \ket{1}\ket{v_j}}.
\end{align}
Equivalently, $F_1$ can be written as 
\begin{align}
    F_1 = Q \Lambda Q^\dagger,
\end{align}
where
\begin{align}
Q &= 
\sum_{j=1}^N \ket{+_j}\bra{0}\bra{j}
+
\sum_{j=1}^N \ket{-_j}\bra{1}\bra{j} \\
&=\frac{1}{\sqrt{2}}\lrb{\ket{0}\bra{0}\otimes U+\ket{1}\bra{0}\otimes W+\ket{0}\bra{1}\otimes U-\ket{1}\bra{1}\otimes W}
\end{align}
is unitary and
\begin{align}
\Lambda &= \lrb{\ket{0}\bra{0} - \ket{1}\bra{1}} \otimes \i 
 \sum_{j=1}^N \gamma_j \ket{j}\bra{j}        
\end{align}
is diagonal. Meanwhile, observe that
\begin{align}
F_2 = \ket{0}\bra{0}\bra{0}\otimes G_1
    +\i\ket{0}\bra{0}\bra{1}\otimes G_2
    +\ket{0}\bra{1}\bra{1}\otimes G_3.    
\end{align}
Thus, we have
\begin{align}
\tilde{F}_2 & = Q^\dagger F_2 Q^{\otimes 2} \\
&=\ket{+}\bra{+}\bra{+} \otimes U^\dagger G_1 (U \otimes U)  +\ket{+}\bra{+}\bra{-} \otimes \i U^\dagger G_2 (U \otimes W) \nonumber\\
&\quad +\ket{+}\bra{-}\bra{-} \otimes U^\dagger G_3 (W \otimes W) \\
& 
=\ket{+}\bra{+}\bra{+} \otimes \tilde{G}_1 +\ket{+}\bra{+}\bra{-} \otimes \i \tilde{G}_2 +\ket{+}\bra{-}\bra{-} \otimes \tilde{G}_3,
\end{align}
where $\ket{\pm}=\frac{1}{\sqrt{2}}(\ket{0} \pm \ket{1})$. Then,
since $\tilde{G}_i$ is $s_i$-column-sparse, for $i=1,2,3$, it follows that $\tilde{F}_2$ is at most $2(s_1+s_2+s_3)$-column-sparse. Furthermore, it is straightforward to verify that
\begin{align}
    \norm{\tilde{F}_2} \le \sqrt{\norm{\tilde{G}_1}^2+\norm{\tilde{G}_2}^2+\norm{\tilde{G}_3}^2}. 
\end{align}
Then we have
\begin{align}
    32 \rho \cdot \operatorname{cs}(\tilde{F}_2) \cdot \norm{\tilde{F}_2} \le c \omega,
\end{align}
as required by Theorem~\ref{thm:algorithm_complexity_time_dependent_f2_case}.
\end{itemize}

To solve Poly-CNOP, we first apply the quantum algorithm from Theorem~\ref{thm:algorithm_complexity_time_dependent_f2_case} to the nonlinear Schr\"{o}dinger equation \eqref{eq:nonlinear_schrodinger_equation_for_nonlinear_oscillators} and obtain a state that is $\frac{b-a}{4}$-close to $\ket{\psi(T)}$. Next, define the projector $\Pi_V= \ket{0}\bra{0}\otimes \sum_{j \in V} \ket{j}\bra{j}$. We then estimate
\begin{align}
   \bra{\psi(T)} \Pi_V \ket{\psi(T)} = \frac{K_V(T)}{E(T)}
\end{align}
within additive error $\frac{b-a}{2}$ by repeatedly performing the projective measurement
$\lrcb{\Pi_V, I-\Pi_V}$ on $\ket{\psi(T)}$ and recording the frequency of observing the outcome corresponding to $\Pi_V$. This projective measurement can be implemented using a single query to the $V$-membership oracle. Moreover, the number of repetitions required for estimating the expectation value can be quadratically reduced using amplitude estimation.

It remains to analyze the cost of the above algorithm. As shown in Section~\ref{subsubsec:coupled_linear_oscillators}, a block-encoding of $F_1=-\i H$ with a rescaling factor $\myO{\mypoly{n}}$ can be implemented using $\myO{1}$ queries to the oracles for $K$ and $M$, along with $\myO{\mypoly{n}}$ elementary gates. 

Meanwhile, a block-encoding of $F_2$ can be constructed from block-encodings of $G_1$, $G_2$ and $G_3$ using standard linear combination of unitaries (LCU) technique. Since each of these block-encodings has a rescaling factor $\myO{\mypoly{n}}$ and can be implemented using $\myO{\mypoly{n}}$ elementary gates, it follows that 
a block-encoding of $F_2$ with a rescaling factor $\myO{\mypoly{n}}$ can also be implemented using $\myO{\mypoly{n}}$ elementary gates. 

Moreover, since $b-a=\myOmega{\frac{1}{\mypoly{n}}}$, it suffices to prepare $\ket{\psi(T)}$ within $\myTheta{\frac{1}{\mypoly{n}}}$ precision and to repeat this procedure $\myO{\mypoly{n}}$ times. 

Putting everything together, and invoking Theorem~\ref{thm:algorithm_complexity_time_dependent_f2_case}, we conclude that the quantum algorithm for Poly-CNOP requires $\myO{\mypoly{n}}$ calls to the oracles for $K$, $M$ and $V$, to the block-encodings of $G_1$, $G_2$ and $G_3$ with $\myO{\mypoly{n}}$ rescaling factors, and to the unitary $\mathcal{W}$ that prepares $\ket{\psi(0)}$, in addition to $\myO{\mypoly{n}}$ elementary gates. Since each of the aforementioned unitaries can be implemented in $\myO{\mypoly{n}}$ time, we conclude that the above algorithm runs in $\myO{\mypoly{n}}$ quantum time. Hence, Poly-CNOP is in BQP. This completes the proof of the theorem.
\end{proof}

\begin{thebibliography}{99}
\bibitem[24]{babbush2023exponential} Ryan Babbush, Dominic W. Berry, Robin Kothari, Rolando D. Somma, and Nathan Wiebe. ``Exponential quantum speedup in simulating coupled classical oscillators''. \href{https://dx.doi.org/10.1103/PhysRevX.13.041041}{Phys. Rev. X {\bf 13}, 041041}~(2023).
\bibitem[44]{wang2017efficient} Guoming Wang. ``Efficient quantum algorithms for analyzing large sparse electrical networks''. \href{https://dx.doi.org/10.26421/QIC17.11-12-5}{Quantum Info. Comput. {\bf 17}, 987--1026}~(2017).
\bibitem[45]{chakraborty2019thepower} Shantanav Chakraborty, Andr\'{a}s Gily\'{e}n, and Stacey Jeffery. ``The power of block-encoded matrix powers: Improved regression techniques via faster {H}amiltonian simulation''. In 46th International Colloquium on Automata, Languages, and Programming (ICALP 2019). \href{https://dx.doi.org/10.4230/LIPIcs.ICALP.2019.33}{Volume 132 of LIPIcs, 33:1--33:14}~(2019).
\end{thebibliography}
\end{document}
